\documentclass[11pt]{article}
\usepackage{fontspec,amsmath,amssymb,amsthm,hyperref}
\usepackage[margin=1in]{geometry}
\setlength{\emergencystretch}{3em}
\title{Characteristic functions are continuous in L\ensuremath{{}^{2}}}
\author{LangevinLean mathematical exposition}
\date{}
\begin{document}\maketitle
\noindent Source module: \texttt{UnadjustedLangevin/LpCharacteristic.lean}. Compile this proof document with XeLaTeX or LuaLaTeX. The accompanying HTML pairs the mathematical argument with the exact Lean statements and records the scope of the formalization.


At each fixed frequency, the characteristic function of the law of a real \(L^2\) variable depends continuously on that variable. Consequently centered Gaussian laws are closed under \(L^2\) convergence, with variance given by the squared norm.


\section{centeredCharacteristic}


Set \(c_t(x)=e^{itx}-1\), represented by the unit-radius circle map centered at \(-1\). Subtracting one makes \(c_t(0)=0\), which is the condition needed by the Lipschitz composition construction on \(L^2\).



\noindent Lean declarations: \texttt{centeredCharacteristic}.

\section{centeredCharacteristic lipschitz}


Multiplication by \(t\) is \(|t|\)-Lipschitz, and the unit circle parametrization is \(1\)-Lipschitz. Composing gives \(|c_t(x)-c_t(y)|\le|t||x-y|\). The proof identifies the nonnegative Lipschitz constants before applying the composition theorem.



\noindent Lean declarations: \texttt{centeredCharacteristic\_lipschitz}.

\section{centeredCharacteristic zero}


At zero the exponential is one, hence \(c_t(0)=e^0-1=0\). Unfolding the circle-map definition makes this a scalar identity; this vanishing proof permits \texttt{compLp} to lift the map without an integrability assumption on a nonzero constant.



\noindent Lean declarations: \texttt{centeredCharacteristic\_zero}.

\section{characteristicLp}


Lift \(c_t\) to \(L^2\) by Lipschitz composition, then add the constant class one. Thus \(\operatorname{characteristicLp}(t,X)\) represents \(e^{itX}\). The probability assumption makes the constant function square integrable.



\noindent Lean declarations: \texttt{characteristicLp}.

\section{characteristicLp coe}


Combine the almost-everywhere representative identities for lifted composition, \(L^2\) addition, and the constant class. Their intersection has full measure; there the representative is \((e^{itX}-1)+1=e^{itX}\). This connects quotient-space operations with the random function used in the characteristic integral.



\noindent Lean declarations: \texttt{characteristicLp\_coe}.

\section{charFun map eq inner characteristicLp}


The pushforward integral theorem gives \(\widehat{P_X}(t)=\int e^{itX}dP\). In complex \(L^2\), inner product with the constant one has exactly this integral because its conjugate is one. The representative identity above therefore yields \(\widehat{P_X}(t)=\langle1,\operatorname{characteristicLp}(t,X)\rangle\).



\noindent Lean declarations: \texttt{charFun\_map\_eq\_inner\_characteristicLp}.

\section{continuous charFun map Lp}


Lipschitz composition is continuous on \(L^2\); adding the constant one preserves continuity. Pairing with one is a continuous linear functional. The preceding integral identity rewrites their composition as \(X\mapsto\widehat{P_X}(t)\), proving continuity of the actual-law characteristic function at every fixed frequency.



\noindent Lean declarations: \texttt{continuous\_charFun\_map\_Lp}.

\section{continuous gaussianReal charFun Lp}


The centered Gaussian characteristic function is \(e^{-t^2\|X\|_2^2/2}\). Norm, squaring, conversion of a nonnegative real to \texttt{\ensuremath{\mathbb R}≥0}, and the exponential are continuous. Lean unfolds the Gaussian formula and composes these continuity facts; the variance is the actual squared \(L^2\) norm.



\noindent Lean declarations: \texttt{continuous\_gaussianReal\_charFun\_Lp}.

\section{hasLaw gaussianReal of tendsto Lp eventually}


At each frequency, \(L^2\) convergence makes the actual characteristic functions converge. It also makes the centered Gaussian characteristic functions with variances \(\|X_i\|_2^2\) converge to variance \(\|Y\|_2^2\). Eventual Gaussianity identifies the two nets eventually, so uniqueness of limits identifies their limits. \texttt{Measure.ext\_of\_charFun} then gives the claimed Gaussian law of \(Y\); a nontrivial filter is required for uniqueness.



\noindent Lean declarations: \texttt{hasLaw\_gaussianReal\_of\_tendsto\_Lp\_eventually}.

\section{hasLaw gaussianReal of tendsto Lp}


If every approximant has the stated Gaussian law, then the law identity holds eventually for any filter. Feed \texttt{Eventually.of\_forall} to the preceding closure theorem. This version is suited to finite orthogonal sums converging to an infinite Gaussian series.



\noindent Lean declarations: \texttt{hasLaw\_gaussianReal\_of\_tendsto\_Lp}.

\end{document}
